\section{Related work}
\label{sec:related}

\paragraph{The systematization this survey builds on}
Papernot, McDaniel, Sinha and Wellman gave the field its standing map of attacks
on machine learning, and later work in the area still uses its vocabulary
\cite{PMSW18}. Their threat model has three parts, the point in the pipeline the
adversary attacks, what it is able to do once there, and the goal it pursues.
The middle part takes a different form in each phase, since their Figure 3
places white box and black box at inference and read and write access at
training. We adopt their attack families wherever they fit, and
Section~\ref{sec:findings-tables} adds one axis they do not carry, the resource
each work bounds and the party it binds.
% claim: resource-axis-absent-from-prior-systematizations

\paragraph{Systematizations of one attack family}
Several later works map a single family in more depth than a general survey
allows, and each reports the same absence from inside its own subfield. Fenaux
and Kerschbaum formalize adversary knowledge for adversarial examples through an
oracle abstraction and a security game, and record a lack of standards for
evaluating those attacks \cite{FK24}. Suya et al.\ place \BBattacks{} black-box attacks
in a four-axis grid and show that several state-of-the-art claims dissolve once
a baseline receives the same knowledge the new attack assumes \cite{SSZHTE24}.
Wen, Liu, Backes and Zhang report that data reconstruction has no agreed formal
definition and supply one with a tolerance \cite{WLBZ25}. Shafran, Shumailov,
Erdogdu and Papernot report the same for model extraction and disentangle the
adversary into prior data, query distribution and query policy \cite{SSEP24}.
Meeus et al.\ do it for membership inference against language models and show
the benchmarks carry a distribution shift a model-free classifier can exploit
\cite{MSJFRM25}. Desfontaines and Pej\'o catalogue the variants of differential
privacy along seven dimensions \cite{DP20}. Each repair supplies knowledge,
access, or a success criterion, and Section~\ref{sec:findings-tables} takes the
budget as its organizing question.

\paragraph{Critiques of evaluation practice}
A second line measures how the field evaluates and not how it defines. Arp et
al.\ name ten pitfalls in learning-based security, measure their prevalence
across \DDpapers{} papers from four venues over a decade, and survey the authors
of those papers for agreement \cite{AQPWPWCR22}. Athalye, Carlini and Wagner
identify obfuscated gradients and circumvent \OGobfuscated{} of the
\OGpopulation{} non-certified defences at one conference \cite{ACW18}. Tram\`er,
Carlini, Brendel and Madry repeat the exercise on \AAdefences{} defences that
had all performed adaptive evaluations, and conclude that the difficulty is
methodological and not technical \cite{TCBM20}. Carlini et al.\ turn the lesson
into a checklist and state a position this survey returns to in
Section~\ref{sec:findings}, that attacker computation should not be a
threat-model parameter \cite{CAPBRTGMK19}. These works diagnose threat models as
unstated or imprecise, and Section~\ref{sec:findings-tables} adds the reading
that a threat model can also be unpriced.

\paragraph{Where the quantity already has a definition}
Fine-grained cryptography is the formal home of the comparison between an
adversary's work and an honest party's. Brzuska and Couteau define an
$(\varepsilon,\delta)$-fine-grained one-way function and then argue at length
that current techniques cannot prove the implication they want \cite{BC22}.
LaVigne, Lincoln and Vassilevska Williams name the gap as an object and build
fine-grained key exchange and public-key encryption around it \cite{LLW19}.
Ball, Rosen, Sabin and Vasudevan supply average-case fine-grained hardness for
the problems those constructions rest on \cite{BRSV17}, and Degwekar,
Vaikuntanathan and Vasudevan build primitives secure against restricted circuit
classes \cite{DVV16}. This literature contains no machine learning, and
Section~\ref{sec:findings-tables} reports that no work in the machine learning
half of the corpus states its adversary in these terms.

\paragraph{Restricted adversary classes as a security tool}
Restricting the adversary to an algorithm class, and not to a running time, has
a record in security. Chaudhuri, Imola and Machanavajjhala bound a privacy
adversary by the function class its attack is drawn from \cite{CIM19}.
Goldwasser, Kim, Vaikuntanathan and Zamir plant backdoors undetectable to any
polynomial-time distinguisher with oracle access \cite{GKVZ22}. Degwekar,
Nakkiran and Vaikuntanathan show that a task where robust classifiers exist but
cannot be learned efficiently implies one-way functions \cite{DNV19}. Jagielski
et al.\ price a model-extraction adversary by algorithm class, showing that
extraction expressible as statistical queries needs exponentially many samples
in depth for fidelity extraction \cite{JCBKP20}, and Chandrasekaran et al.\ give
extraction a query budget inside its definition by reducing it to
query-synthesis active learning \cite{CCGJY20}. Dziedzic et al.\ charge an
extraction adversary compute directly, calibrating a proof-of-work puzzle to how
much information a client has already extracted \cite{DKLP22}. These are the
works closest to the question this survey asks, and
Section~\ref{sec:findings-tables} places each of them alongside the rest of the
corpus, so a reader can compare their adversary assumptions with those of works
that state none.
